\section{模仿学习}

\subsection{从试错到示教}

\begin{figure}[H]
\centering
\includegraphics[width=0.95\textwidth]{slides-images/slide-046.jpg}
\caption{模仿学习的基本思想：从人类示教数据中学习策略 $\pi_\theta: o \mapsto a$\protect\footnotemark}
\end{figure}
\footnotetext{Slides 第47页。}

模仿学习（Imitation Learning）是当前机器人学习中\textbf{最实用}的方法。其核心思想：收集大量人类执行任务的示教数据 $(o_t, a_t)$，用监督学习训练策略网络 $\pi_\theta$，直接从观测映射到动作。

\begin{importantbox}{行为克隆（Behavior Cloning）}
最简单的模仿学习算法。给定专家示教数据集 $\{(o_i, a_i)\}$，通过最小化预测动作与专家动作之间的损失来训练策略：
$$\min_\theta \sum_i \mathcal{L}(\pi_\theta(o_i), a_i)$$
优点是简单高效——早上收集数据，下午训练模型，晚上就有工作策略。
\end{importantbox}

\subsection{级联误差问题}

\begin{figure}[H]
\centering
\includegraphics[width=0.95\textwidth]{slides-images/slide-048.jpg}
\caption{级联误差（Cascading Error）：小误差逐步累积放大，导致策略偏离示教轨迹\protect\footnotemark}
\end{figure}
\footnotetext{Slides 第49页。}

\begin{warningbox}{行为克隆的致命缺陷：级联误差}
由于机器人学习是序列决策问题，一个微小的预测误差会导致状态偏离训练数据的分布，在新状态下策略产生更大的误差，进而更大偏离——误差沿时间轴指数级放大。这与标准监督学习中数据点独立的假设截然不同。
\end{warningbox}

解决方案是\textbf{迭代式数据收集}：

\begin{enumerate}
    \item 收集初始专家示教
    \item 训练策略并在真实环境中执行
    \item 观察失败案例，收集纠正行为数据
    \item 将纠正数据加入训练集，重新训练
    \item 重复上述过程直到策略足够鲁棒
\end{enumerate}

\subsection{逆强化学习}

除了直接学习动作映射，另一条思路是\textbf{逆强化学习（Inverse RL）}：从示教数据中推断隐含的奖励函数，再用标准强化学习来学习策略。经典案例包括 Stanford 的 Pieter Abbeel 和 Andrew Ng 利用此方法控制遥控直升机完成极限飞行动作。

\subsection{扩散策略（Diffusion Policy）}

\begin{figure}[H]
\centering
\includegraphics[width=0.95\textwidth]{slides-images/slide-054.jpg}
\caption{扩散策略：借鉴扩散模型的去噪过程来生成动作序列，显著提升了策略对多模态示教的处理能力\protect\footnotemark}
\end{figure}
\footnotetext{Slides 第55页。}

近年来模仿学习最重要的进展之一是\textbf{扩散策略}（Diffusion Policy），将扩散模型引入策略函数类：

\begin{knowledgebox}{从生成模型到策略学习}
\begin{itemize}
    \item \textbf{隐式行为克隆}（Implicit BC）：借鉴能量基模型（EBM），学习 $(o, a) \mapsto \text{score}$，通过优化推理动作
    \item \textbf{扩散策略}：借鉴扩散模型，通过逐步去噪生成动作序列，天然处理多模态分布
\end{itemize}
两者都从深度学习的生成模型发展中汲取灵感，显著提升了对复杂操作任务的处理能力。
\end{knowledgebox}

扩散策略展示了极为丰富的操作能力：涂抹黄油、煎蛋、削土豆皮、滑动书本等精细操作，都能在真实世界中可靠执行。

\subsection{本章小结}

模仿学习是目前获得真实世界工作策略的最高效途径。行为克隆虽然简单，但面临级联误差问题，需要迭代式数据收集来缓解。扩散策略等新方法从生成模型领域引入先进技术，显著提升了策略的表达能力和鲁棒性。

%% ========== Part 6: 机器人基础模型 ==========

